\section{Descomposición del ELBO y objetivo de entrenamiento}

\subsection{De la verosimilitud logarítmica al ELBO}

Similar a los VAE, el objetivo de entrenamiento de los modelos difusos es maximizar la verosimilitud logarítmica de los datos $\log p_\theta(x_0)$. Dado que optimizarlo directamente no es factible, pasamos a optimizar su cota inferior: la cota de evidencia (ELBO):

$$\log p_\theta(x_0) \geq \text{ELBO} = \mathbb{E}_{q(x_{1:T}|x_0)}\left[\log \frac{p_\theta(x_{0:T})}{q(x_{1:T}|x_0)}\right]$$

Equivalentemente, durante el entrenamiento minimizamos el ELBO negativo:

$$L = -\text{ELBO}$$

\begin{knowledgebox}{Relación entre ELBO y VAE}
La derivación del ELBO es idéntica a la de los VAE: se introduce una posterior variacional $q(x_{1:T}|x_0)$ y se utiliza la desigualdad de Jensen para obtener la cota inferior. La diferencia radica en que, para los modelos difusos, $q$ es un proceso hacia adelante predefinido (no aprendido), por lo que el gap del ELBO está determinado por el diseño del proceso hacia adelante.
\end{knowledgebox}

\subsection{Descomposición término a término del ELBO}

Aprovechando la propiedad de Markov y las reglas de Bayes, el ELBO puede descomponerse en la suma de $T+1$ términos. Concretamente, el ELBO negativo se puede expresar como:

$$L = \underbrace{D_{\text{KL}}\big(q(x_T|x_0) \,\|\, p(x_T)\big)}_{L_T\text{（término de ajuste al prior）}} + \sum_{t=2}^{T} \underbrace{D_{\text{KL}}\big(q(x_{t-1}|x_t, x_0) \,\|\, p_\theta(x_{t-1}|x_t)\big)}_{L_{t-1}\text{（término de ajuste al desruido）}} + \underbrace{\big(-\log p_\theta(x_0|x_1)\big)}_{L_0\text{（término de reconstrucción）}}$$

Significado de cada término:
\begin{itemize}
    \item $L_T$ (término de ajuste al prior): divergencia KL entre la distribución terminal del proceso hacia adelante $q(x_T|x_0)$ y el prior $p(x_T) = \mathcal{N}(0,I)$. Dado que $q$ está predefinido y cuando $T$ es suficientemente grande $q(x_T|x_0) \approx \mathcal{N}(0,I)$, este término se aproxima a 0 y \textbf{no requiere optimización}.
    \item $L_{t-1}$ (término de ajuste al desruido): este es el término central; exige que la transferencia inversa aprendida $p_\theta(x_{t-1}|x_t)$ coincida con la transferencia posterior real $q(x_{t-1}|x_t, x_0)$.
    \item $L_0$ (término de reconstrucción): verosimilitud logarítmica negativa del último paso para reconstruir $x_0$ a partir de $x_1$.
\end{itemize}

\begin{importantbox}{El núcleo del entrenamiento: el término de ajuste al desruido $L_{t-1}$}
La meta clave del entrenamiento es minimizar el término de ajuste al desruido $L_{t-1}$, haciendo que la distribución de transferencia inversa $p_\theta(x_{t-1}|x_t)$ se aproxime lo más posible a la transferencia posterior real $q(x_{t-1}|x_t, x_0)$. Esto requiere primero derivar una expresión cerrada para $q(x_{t-1}|x_t, x_0)$.
\end{importantbox}

\subsection{Derivación de la distribución posterior real de transferencia}

Utilizando las reglas de Bayes:

$$q(x_{t-1}|x_t, x_0) = \frac{q(x_t|x_{t-1}, x_0) \cdot q(x_{t-1}|x_0)}{q(x_t|x_0)}$$

Dado que las tres distribuciones $q(x_t|x_{t-1})$, $q(x_{t-1}|x_0)$ y $q(x_t|x_0)$ son distribuciones gaussianas, aprovechando las reglas de multiplicación y división de distribuciones gaussianas, se puede demostrar que $q(x_{t-1}|x_t, x_0)$ también es una distribución gaussiana:

$$q(x_{t-1}|x_t, x_0) = \mathcal{N}\big(x_{t-1};\, \tilde{\mu}_t(x_t, x_0),\; \tilde{\beta}_t I\big)$$

Donde la varianza y la media son respectivamente:

$$\tilde{\beta}_t = \frac{(1 - \bar{\alpha}_{t-1}) \cdot \beta_t}{1 - \bar{\alpha}_t}$$

$$\tilde{\mu}_t(x_t, x_0) = \frac{\sqrt{\bar{\alpha}_{t-1}}\, \beta_t}{1 - \bar{\alpha}_t}\, x_0 + \frac{\sqrt{\alpha_t}\, (1 - \bar{\alpha}_{t-1})}{1 - \bar{\alpha}_t}\, x_t$$

\begin{warningbox}{Dependencia posterior de $x_0$}
Obsérvese que la media $\tilde{\mu}_t$ de $q(x_{t-1}|x_t, x_0)$ depende tanto de $x_t$ como de $x_0$. Sin embargo, en el proceso inverso (proceso generativo), \textbf{no} tenemos información sobre $x_0$; este es precisamente el motivo por el cual necesitamos una red neuronal para predecirlo. La varianza $\tilde{\beta}_t$ solo depende de los parámetros del esquema de ruido, no de $x_0$, por lo que no requiere aprendizaje.
\end{warningbox}

\subsection{Resumen de la sección}

La descomposición del ELBO transforma el objetivo de entrenamiento en minimizar la divergencia KL entre la distribución de transferencia inversa y la distribución posterior real de transferencia en cada paso temporal. La transferencia posterior real $q(x_{t-1}|x_t, x_0)$ es una distribución gaussiana cuya varianza se puede calcular directamente y cuya media es una combinación lineal de $x_t$ y $x_0$.

%% ========== Sección 4: Parametrización del proceso inverso ==========


